\begin{tikzpicture}[
  font=\scriptsize,
  execute at begin node={\hyphenpenalty=10000\exhyphenpenalty=10000},
  aktor/.style={draw, rounded corners, fill=blue!8, minimum width=1.6cm, minimum height=0.6cm, align=center},
  judul/.style={font=\scriptsize\bfseries, align=center},
  garis/.style={dashed, gray},
  ket/.style={font=\tiny, align=center},
  panah/.style={-{Stealth[length=2mm]}, thick},
  balik/.style={-{Stealth[length=2mm]}, thick, dashed}
]
  % Panel kiri: pemanggil menunggu sepanjang penyedia bekerja.
  \node[judul] at (2.6,0.75) {Sinkron: pemanggil menunggu jawaban};
  \node[aktor] (pa) at (0.8,0) {Pemanggil};
  \node[aktor] (sa) at (4.4,0) {Penyedia};
  \draw[garis] (pa.south) -- (0.8,-4.3);
  \draw[garis] (sa.south) -- (4.4,-4.3);
  \draw[fill=red!22] (0.68,-0.8) rectangle (0.92,-3.3);
  \draw[fill=orange!12] (4.28,-0.8) rectangle (4.52,-3.3);
  \draw[panah] (0.92,-0.8) -- node[above, ket] {\textit{request}} (4.28,-0.8);
  \node[ket, anchor=west] at (1.0,-1.7) {pemanggil tertahan};
  \node[ket, anchor=east] at (4.2,-2.4) {penyedia bekerja};
  \draw[balik] (4.28,-3.3) -- node[above, ket] {\textit{response}: 200 dan hasil} (0.92,-3.3);
  \node[ket] at (2.6,-3.85) {Hasil tiba dalam satu interaksi};
  \node[ket] at (2.6,-4.6) {Bila penyedia lambat, pemanggil ikut lambat};

  % Panel kanan: penerimaan dan penyelesaian terjadi pada waktu berbeda.
  \begin{scope}[xshift=6.4cm]
    \node[judul] at (2.6,0.75) {Asinkron: diterima dulu, selesai belakangan};
    \node[aktor] (pb) at (0.8,0) {Pemanggil};
    \node[aktor] (sb) at (4.4,0) {Penyedia};
    \draw[garis] (pb.south) -- (0.8,-4.3);
    \draw[garis] (sb.south) -- (4.4,-4.3);
    \draw[fill=red!22] (0.68,-0.8) rectangle (0.92,-1.3);
    \draw[fill=orange!12] (4.28,-0.8) rectangle (4.52,-3.6);
    \draw[panah] (0.92,-0.8) -- node[above, ket] {\textit{request}} (4.28,-0.8);
    \draw[balik] (4.28,-1.3) -- node[above, ket] {202 dan ID pekerjaan} (0.92,-1.3);
    \node[ket, anchor=west] at (1.0,-1.75) {pemanggil mengerjakan hal lain};
    \draw[panah] (0.8,-2.45) -- node[above, ket] {tanya status memakai ID} (4.28,-2.45);
    \draw[balik] (4.28,-2.85) -- node[above, ket] {masih diproses} (0.8,-2.85);
    \draw[balik, draw=praditagreen] (4.28,-3.6) -- node[above, ket] {selesai atau gagal} (0.8,-3.6);
    \node[ket] at (2.6,-4.6) {Jawaban pertama hanya berarti ``diterima'', bukan ``selesai''};
  \end{scope}
\end{tikzpicture}
